% 用 xelatex 编译
%\documentclass[a4paper,cs4size,hyperref,fancyhdr]{ctexbook}
\documentclass[UTF8,openright]{ctexbook}

% 论文版面要求:
% 统一按 word 格式A4纸（页面设置按word默认值）编排、打印、制作。
% 正文内容字体为宋体；字号为小4号；字符间距为标准；行距为25磅(约0.88175cm)。

%%%%% ===== 页面设置
\usepackage[a4paper,top=2.54cm,bottom=2.54cm,left=3.17cm,right=3.17cm,%
            ]{geometry}
            
\setlength{\parindent}{2em}
%默认的弹性间距会导致文中某些排版flush的时候，出现大量空白。
\setlength{\parskip}{0.5em} %指定固定段后间距，默认为弹性间距。
\setlength{\intextsep}{10pt} %固定浮浮动体前后间距。


%%%%% =====章节 标题 设置
\ctexset{%
  contentsname={\vspace{-3.5em}\centerline{\zihao{-3}\heiti 目\quad 录}\vspace{-0.7em}},
  listfigurename={\vspace{-3.5em}\centerline{\zihao{-3}\heiti 插\ 图\ 目\ 录}\vspace{-0.5em}},
  listtablename={\vspace{-3.5em}\centerline{\zihao{-3}\heiti 表\ 格\ 目\ 录}\vspace{-0.5em}},
  bibname={\vspace{-3em}\centerline{\zihao{-3}\heiti 参\ 考\ 文\ 献}\vspace{3em}},
  chapter={name={第,章},
           number=\chinese{chapter}, %指定章序号为一二三。。。。
           nameformat={\zihao{3}\bfseries},
           titleformat={\zihao{3}\bfseries},
           beforeskip={-10pt},
           afterskip={20pt}
           },
  section={format=\raggedright,
           nameformat={\zihao{4}\bfseries},
           titleformat={\zihao{4}\bfseries},
%           afterskip={1ex plus 0.2ex}
           beforeskip={1ex},% 固定段前段后间距，
           afterskip={1ex}
           },
  subsection={format=\raggedright,
           nameformat={\zihao{-4}\bfseries},
           titleformat={\zihao{-4}\bfseries},
%           afterskip={0.5ex plus 0.1ex}
           beforeskip={0.5ex},
           afterskip={0.5ex}
           }
}
%%%%% ===== 中英文字体
\setmainfont{Times New Roman}
%\setsansfont{Myriad Pro} % 无衬线字体 sans serif \sffamily
%\setmonofont{Consolas}   % 等宽字体 typewriter \ttfamily
%\newcommand{\Times}{\fontspec{Times New Roman}}
%% 中文字体
%\setCJKmainfont[BoldFont={Microsoft YaHei},ItalicFont={KaiTi}]{NSimSun}
%\setCJKsansfont{Microsoft YaHei}
%\setCJKmonofont{KaiTi}
%\setCJKfamilyfont{STSong}{方正小标宋_GBK}\newcommand{\STSong}{\CJKfamily{STSong}}
\setCJKfamilyfont{STSong}{STZhongsong}\newcommand{\STSong}{\CJKfamily{STSong}}

%%%%% ===== 常用宏包
\usepackage{amsmath,amssymb,amsfonts,bm}
\usepackage[amsmath,thref,thmmarks,hyperref]{ntheorem}
\usepackage{graphicx,xcolor,float}
\usepackage{fancyhdr}
\usepackage{tocloft} % 设置目录中的条目间距


\renewcommand\cftdot{\textsubscript{……}}
\renewcommand\cftdotsep{0}

\setlength{\cftbeforechapskip}{1pt}
\renewcommand{\cftchapleader}{\cftdotfill{\cdot}}


\usepackage{booktabs} % toprule, midrule, bottomrule
\usepackage{varwidth} % 可变宽度的 parbox

%%%%% ===== 参考文献与链接
\usepackage[numbers,sort&compress,sectionbib,super, square]{natbib} %引用上标，禁用连续缩写。
\newcommand{\upcite}[1]{\textsuperscript{\cite{#1}}}


\usepackage[xetex,pagebackref]{hyperref}
\hypersetup{CJKbookmarks=true,colorlinks=true,citecolor=blue,%
            linkcolor=blue,urlcolor=blue,bookmarksnumbered=true,%
	        bookmarksopen=true,breaklinks=true}
	        
	        
	        
\iffalse   % 调试时，可去掉，以用于显示引用位置。
\renewcommand*{\backrefalt}[4]{%
\ifcase #1 No citations.%
\or Cited on page #2.%
\else Cited on pages #2.%
%\else #1 Cited on pages #2.%
\fi
}

\else
\renewcommand*{\backrefalt}[4]{}
\fi

%%%%% ===== 浮动图表的标题
\usepackage[margin=2em,labelsep=space,skip=0.5em,font=normalfont]{caption}
\DeclareCaptionFormat{mycaption}{{\heiti\color{blue} #1}#2{\kaishu #3}}
\captionsetup{format=mycaption,tablewithin=chapter,figurewithin=chapter}%,belowskip=-10pt
%\setlength{\belowcaptionskip}{-10pt}

%%%%%% ===== 浮动图表的比例默认50%以下，否则无法浮动。
\renewcommand\floatpagefraction{.9} %当浮动体小于页面90%时进行直接放置。
\renewcommand\topfraction{.9}  
\renewcommand\bottomfraction{.9}  
\renewcommand\textfraction{.1}








%%%%% ===== 算法
\usepackage{algorithm,algpseudocode}

%%%%% ===== 其他
\usepackage{ulem}
\def\ULthickness{1pt}




%%%%%===== Code Style代码
\usepackage{listings}
\usepackage{color}

\definecolor{dkgreen}{rgb}{0,0.6,0}
\definecolor{gray}{rgb}{0.5,0.5,0.5}
\definecolor{mauve}{rgb}{0.58,0,0.82}

\lstset{frame=tb,
  language=Scala,
  aboveskip=3mm,
  belowskip=3mm,
  showstringspaces=false,
  columns=flexible,
  basicstyle={\small\ttfamily},
  numbers=none,
  numberstyle=\tiny\color{gray},
  keywordstyle=\color{blue},
  commentstyle=\color{dkgreen},
  stringstyle=\color{mauve},
  breaklines=true,
  breakatwhitespace=true,
  tabsize=3
}



%%%%% ===== 自定义命令
%\renewcommand{\C}{\mathbb{C}}
%\newcommand{\R}{\mathbb{R}}


%%%%%%%%%%%%%%%%%%%%%%%%%%%%%%%%%%%%%%%%%%%%%%%%%%%%%%%%%%%%%%
\begin{document}

%%%%% ===== 格式
%%%%% ===== 设置页面

%%%%% ===== 设置数学公式
\numberwithin{equation}{chapter}
\allowdisplaybreaks




%%%%% ===== 设置页眉和页脚
\pagestyle{fancy}
\fancyhf{}  % 清除以前对页眉页脚的设置

\newcommand{\makeheadrule}{%% 定义页眉与正文间双隔线
  \makebox[0pt][l]{\rule[.7\baselineskip]{\headwidth}{0.3pt}}%0.4
  \rule[0.85\baselineskip]{\headwidth}{1.0pt}\vskip-.8\baselineskip}
\makeatletter
\renewcommand{\headrule}{%
  {\if@fancyplain\let\headrulewidth\plainheadrulewidth\fi\makeheadrule}}
\makeatother
\renewcommand{\chaptermark}[1]{\markboth{\CTEXthechapter \ #1}{}}
\renewcommand{\sectionmark}[1]{\markright{\thesection \ #1}{}}
%\fancyhead[RO,LE]{{\small\songti\rightmark}}     % 节标题
%\fancyhead[RE]{{\small\songti\leftmark}}      % 章标题
\fancyhead[CO,CE]{华东师范大学硕士专业学位论文}
%\fancyhead[RO,LE]{$\cdot$ {\small\thepage} $\cdot$}
\fancyfoot[RO,LE]{{\thepage}}
%\fancyfoot[CO,CE]{{\thepage}}



%%%%% ===== 定理类环境
\theorempostskipamount0em
\theoremstyle{plain}
\theoremheaderfont{\normalfont\heiti\color{blue}}
\theorembodyfont{\normalfont\kaishu\color{black}}
\theoremindent0em
\theoremseparator{\hspace{0.2em}}
\theoremnumbering{arabic}
\newtheorem{theorem}{定理}[chapter]
\newtheorem{lemma}[theorem]{引理}
\newtheorem{corollary}[theorem]{推论}
\newtheorem{proposition}[theorem]{命题}
\newtheorem{property}[theorem]{性质}
\newtheorem{definition}{定义}[chapter]
\newtheorem{remark}{注记}[chapter]
%
%\theoremheaderfont{\normalfont\itshape\color{blue}}
\theorembodyfont{\normalfont\rmfamily\color{black}}
\newtheorem{example}{例}[chapter]

\theoremstyle{nonumberplain}
\theorempreskip{0em}
\theoremsymbol{\ensuremath{\Box}}
\newtheorem{proof}{证明}

%%%%% ===== 算法
\floatname{algorithm}{\color{blue} 伪代码}
\algrenewcommand{\algorithmiccomment}[1]{\quad{\color{red}\%\ #1}}
\numberwithin{algorithm}{chapter}
\renewcommand{\listalgorithmname}{算\ 法\ 目\ 录}


\def\tableautorefname{表}%
\def\figureautorefname{图}%

\def\equationautorefname{式}%
%\def\footnoteautorefname{脚注}%
%\def\itemautorefname{项}%
%\def\figureautorefname{图}%
%\def\tableautorefname{表}%
%\def\partautorefname{篇}%
%\def\appendixautorefname{附录}%
\def\chapterautorefname{章}%
\def\sectionautorefname{节}%
\def\subsectionautorefname{小节}%
%\def\subsubsectionautorefname{小小节}%
%\def\paragraphautorefname{段落}%
%\def\subparagraphautorefname{子段落}%
%\def\FancyVerbLineautorefname{行}%
%\def\theoremautorefname{定理}%
\def\algorithmautorefname{伪代码}

%%%%%%%%%%%%%%%%%%%%%%%%%%%%%%%%%%%%%%%%%%%%%%%%%%%%%%%%%%%%%%%%%%%%%%%
%%%%% ===== 封面 =====
\makeatletter
\def\class#1{\def\@class{#1}}\def\@class{}
\def\security#1{\def\@security{#1}}\def\@security{}
\def\graduateyear#1{\def\@graduateyear{#1}}\def\@graduateyear{}
\def\ctitle#1{\def\@ctitle{#1}}\def\@ctitle{}
\def\caffil#1{\def\@caffil{#1}}\def\@caffil{}
\def\cmajor#1{\def\@cmajor{#1}}\def\@cmajor{}
\def\cdirection#1{\def\@cdirection{#1}}\def\@cdirection{}
\def\cauthor#1{\def\@cauthor{#1}}\def\@cauthor{}
\def\studentid#1{\def\@studentid{#1}}\def\@studentid{}
\def\csupervisor#1{\def\@csupervisor{#1}}\def\@csupervisor{}
\def\cdate#1{\def\@cdate{#1}}\def\@cdate{}
\def\edegree#1{\def\@edegree{#1}}\def\@edegree{}
\def\etitle#1{\def\@etitle{#1}}\def\@etitle{}
\def\eaffil#1{\def\@eaffil{#1}}\def\@eaffil{}
\def\emajor#1{\def\@emajor{#1}}\def\@emajor{}
\def\edirection#1{\def\@edirection{#1}}\def\@edirection{}
\def\eauthor#1{\def\@eauthor{#1}}\def\@eauthor{}
\def\esupervisor#1{\def\@esupervisor{#1}}\def\@esupervisor{}
\def\edate#1{\def\@edate{#1}}\def\@edate{}

\newcommand{\mcc}[1]{\multicolumn{1}{c}{\underline{\makebox[10em][c]{#1}}}}
\newcommand{\mce}[1]{\multicolumn{1}{c}{\underline{\makebox[15em][l]{#1}}}}

\def\makecover{
\begin{titlepage}

%%%%% ===== 中文封面
  \pdfbookmark[0]{中文封面}{ccover}
  \linespread{1.1}\zihao{4}\ziju{0.05}
  \noindent{\@graduateyear\,届研究生硕士学位论文}\smallskip
  \par\noindent\hspace*{-12pt}
  \setlength{\tabcolsep}{2pt}
  \begin{tabular}[t]{rl}
    分\hspace{7.5pt}类\hspace{7.5pt}号：
                      &\ \underline{\makebox[6em][l]{\,\@class}}  \\[1.5ex]
    密\hspace{2em}级：&\ \underline{\makebox[6em][l]{\@security}}
  \end{tabular}
  \hfill
  \begin{tabular}[t]{rl}
    学校代码：&\ \underline{\makebox[6em][l]{\,10269}} \\[1.5ex]
    学\hspace{2em}号：&\ \underline{\makebox[6em][l]{\,\@studentid}}
  \end{tabular}

  \vspace{4em}
  \begin{center}
    \raisebox{1ex}[0pt][0pt]{
    \includegraphics[width=0.15\textwidth]{ecnu_logo}}\ \
    \includegraphics[width=0.75\textwidth,height=5em]{ecnu}

    \bigskip
    {\STSong
      \zihao{-2}\textbf{East China Normal University} \\[0.5ex]
      \zihao{2} 硕士学位论文\\[1ex]
      \zihao{-2} \textbf{MASTER'S  DISSERTATION}

    % title
    \vspace{3em}\noindent
    \parbox[t]{0.25\textwidth}{\zihao{2} 论文题目：}
    \begin{varwidth}[t]{.75\linewidth}\linespread{1.3}\zihao{2} \@ctitle \linebreak\end{varwidth}


    \vfill\linespread{1.5}\selectfont\zihao{4}
    \renewcommand{\arraystretch}{1.2}
    \begin{tabular}{p{0cm}p{5.5em}@{\extracolsep{0.5ex}}cc}
     ~ & 院 \hfill 系： & & \mcc{\@caffil } \\
     ~ & 专 \hfill 业： & & \mcc{\@cmajor}\\
     ~ & 研 \hfill 究 \hfill 方 \hfill 向：& & \mcc{\@cdirection}\\
     ~ & 指\hfill 导\hfill 老\hfill 师：& & \mcc{\@csupervisor}\\
     ~ & 学\hfill 位\hfill 申\hfill 请\hfill 人：& & \mcc{\@cauthor}\\
    \end{tabular}
    }

    \vfill\@cdate
  \end{center}



















\clearpage % {\pagestyle{empty}\cleardoublepage}
%%%%% ===== 英文封面
  \thispagestyle{empty}
  \pdfbookmark[0]{英文封面}{ecover}
  \noindent Dissertation for master degree in \@graduateyear
  \hfill University Code:\, 10269\par\medskip
  \mbox{}\hfill Student ID:\, \@studentid


  \vspace{5em}
  \begin{center}
    {\zihao{-0} East China Normal University}
    %\scalebox{1.15}[1.6]{\bfseries EAST CHINA NORMAL UNIVERSITY}

    % title
    \vspace{6em}{\bfseries
    \parbox[t]{0.125\textwidth}{\zihao{2} Title：}
    \begin{varwidth}[t]{.75\linewidth}\linespread{1.4}\zihao{2} \@etitle \linebreak\end{varwidth}}

    \vfill\linespread{1.5}\selectfont\mdseries

    \renewcommand{\arraystretch}{1.1}
    \scalebox{1}[1.0]{\setlength{\tabcolsep}{0.5ex}
    \begin{tabular}{rl}
     Department: &  \mce{\@eaffil } \\
     Major:      &  \mce{\@emajor}\\
     Research Direction: &  \mce{\@edirection}\\
     Supervisor: &  \mce{\@esupervisor}\\
     Candidate:  & \mce{\@eauthor}\\
    \end{tabular}
    }

    \vfill
    \@edate
  \end{center}

\end{titlepage}
\makeatother
}

\zihao{-4}

%%%%% ===== 中文封面信息
\graduateyear{2023} % 毕业年份
\class{   } % 分类号（TPxxx.x）
\ctitle{\uline{标题标题标题标题标题标题\linebreak 标题标题标题标题标题标题}}
\def\cctitle{标题标题标题标题标题标题标题标题标题标题标题标题} % 在原创性声明中使用, 不能出现手工换行
\caffil{数学科学学院}
\cmajor{基础数学} % 计算数学
\cdirection{分形几何及其应用} % 数值代数
\csupervisor{李文侠\, 教授}
\cauthor{梁越}
\def\ccauthor{\@cauthor} % 在答辩委员会中使用
\studentid{51205500020}
\cdate{2023 年 5 月}

%%%%% ===== 英文封面信息
\etitle{\uline{Research and Implementation of  \linebreak What You want to do}}
\eaffil{School of Mathematical Sciences} 
\emajor{Fundamental Mathematics} % Computational Mathematics
\edirection{Fractal Geometry and Its Applications} % Numerical Algebra
\esupervisor{Wenxia Li (Professor)}
\eauthor{Yue Liang}
\edate{May, 2023}

%%%%% ===== 生成封面 ===== %使用不同的页边距
\newgeometry{top=2.0cm,bottom=2.0cm,left=2.5cm,right=2.5cm}
{
\renewcommand{\baselinestretch}{1.6}
\makecover
}

\restoregeometry %还原页边距





%%%%% ===== 原创性声明与著作权使用声明 =====
%%%%% ===== 原创性声明与著作权使用声明
\thispagestyle{empty}
\pdfbookmark[0]{原创性声明}{Declaration}

\vspace*{1em}

{
\linespread{1.4}\zihao{-4}
\centerline{\zihao{3}\STSong 华东师范大学学位论文原创性声明}

\bigskip

郑重声明：本人呈交的学位论文《\cctitle》，
是在华东师范大学攻读硕士/博士（请勾选）
学位期间，在导师的指导下进行的研究工作及取得的研究成果。
除文中已经注明引用的内容外，本论文不包含其他个人已经发表或撰写过的研究成果。
对本文的研究做出重要贡献的个人和集体，均已在文中作了明确说明并表示谢意。

\vspace{1em}


{\STSong 作者签名}：$\underline{\hspace{4cm}}$ \hfill
{\STSong 日\quad 期}： \qquad\ 年 \quad\ 月 \quad\ 日 \qquad\mbox{}

\vspace{4em}

\centerline{\zihao{3}\STSong 华东师范大学学位论文著作权使用声明}
\bigskip


《\cctitle》
系本人在华东师范大学攻读学位期间在导师指导下完成的硕士/博士（请勾选）学位论文，
本论文的著作权归本人所有。
本人同意华东师范大学根据相关规定保留和使用此学位论文，
并向主管部门和学校指定的相关机构送交学位论文的印刷版和电子版；
允许学位论文进入华东师范大学图书馆及数据库被查阅、借阅；
同意学校将学位论文加入全国博士、硕士学位论文共建单位数据库进行检索，
将学位论文的标题和摘要汇编出版，采用影印、缩印或者其它方式合理复制学位论文。

本学位论文属于（请勾选）
\newcounter{muni}
\begin{list}{{\hfill\upshape (\qquad)\ \arabic{muni}. }}{%
     \usecounter{muni}\leftmargin6.5em\labelwidth4.2em\labelsep0.2em
     \itemsep0.5em\itemindent0pt\parsep0pt\topsep0pt}

\item 经华东师范大学相关部门审查核定的“内部”或“涉密”学位论文*，
  于 \qquad 年 \quad 月 \quad  日解密, 解密后适用上述授权。

\item 不保密，适用上述授权。
\end{list}

\vskip0.8cm


{\STSong 导师签名}：$\underline{\hspace{4cm}}$ \hfill
{\STSong 本人签名}：$\underline{\hspace{4cm}}$

\bigskip

{\mbox{}\hfill 年\qquad 月\qquad  日 }

\vfill

\parbox[t]{0.946\textwidth}{\zihao{5}
*“涉密”学位论文应是已经华东师范大学学位评定委员会办公室或保密委员会
审定过的学位论文（需附获批的《华东师范大学研究生申请学位论文“涉密”审批表》
方为有效），未经上述部门审定的学位论文均为公开学位论文。
此声明栏不填写的，默认为公开学位论文，均适用上述授权）。\\
}}

\clearpage{\pagestyle{empty}\cleardoublepage}

%%%%% ===== 答辩委员会成员 =====
%%%% ===== 答辩委员会成员
\thispagestyle{empty}
\pdfbookmark[0]{答辩委员会}{Committee}
\vspace*{2em}

\makeatletter
\begin{center}\STSong\zihao{3}
 \underline{\ \ccauthor\ } 硕士学位论文答辩委员会成员名单
\end{center}
\makeatother

\begin{center}\zihao{4}
\renewcommand{\arraystretch}{1.4}
  \begin{tabular}{|c|c|c|c|} \hline
   ~~~~~姓~名~~~~~ & ~~~~~职~称~~~~~ 
   & \hspace{6em}单~位\hspace{6em} & ~~~备~注~~~\\\hline
         XXX    &   教授    &  XXXXX大学数学系  & 主席  \\ \hline
         XXX    &   教授    &  XXXXX大学数学系  &       \\ \hline
         XXX    &   教授    &  XXXXX大学数学系  &       \\ \hline
                &           &                   &       \\ \hline
                &           &                   &       \\ \hline
  \end{tabular}
\end{center}

\clearpage{\pagestyle{empty}\cleardoublepage}


\frontmatter

%\restoregeometry

\makeatletter
\let\ps@plain\ps@fancy% Let 'plain' be exactly the same as 'fancy'
\makeatother

\pagestyle{fancy} %使用fancy页眉和页脚
\fancyfoot[LO,LE]{}
\fancyfoot[RO,RE]{{\thepage}} %单页打印，此处页码固定右下角



%%%%% ===== 中文摘要 =====

\include{Abstract_chs}
\clearpage{\pagestyle{empty}\cleardoublepage}


%%%%% ===== 英文摘要 =====
\include{Abstract_eng}
\clearpage{\pagestyle{empty}\cleardoublepage}













%%%%% ===== 生成目录

\let\cleardoublepage\clearpage

%\pagestyle{empty}
\setcounter{tocdepth}{2} %向下显示到两级标题
\phantomsection % 修复引用问题


\addcontentsline{toc}{chapter}{目录}

\tableofcontents

\newpage%{\pagestyle{empty}\cleardoublepage}
\pagestyle{fancy}



























%\clearpage%{\pagestyle{empty}\cleardoublepage}

%%%%%% ===== 正文部分 ===== %%%%%
\mainmatter
\linespread{1.6}\selectfont % 行距固定25磅，宋体
%\setlength{\baselineskip}{0.88175cm}
\pagestyle{fancy}
\fancyfoot[RO,RE]{}
\fancyfoot[LO,LE]{} %清空之前的页码排版。
\fancyfoot[RO,LE]{{\thepage}}%双面打印，页码向外排布。

\chapter{引言}
\section{研究背景和意义}


引言部分\cite{bi:hw-ml-challenges}, 介绍论文研究课题的应用背景或者问题来源,\cite{bi:hw-sw-co-design}
一些基本概念, 现有成果等等.


\section{国内外研究现状}


现有成果介绍, 现有成果介绍, 现有成果介绍, 现有成果介绍, 现有成果介绍.
现有成果介绍, 现有成果介绍, 现有成果介绍, 现有成果介绍, 现有成果介绍.
现有成果介绍, 现有成果介绍, 现有成果介绍, 现有成果介绍, 现有成果介绍.
\begin{theorem}
  这是定理，这是定理，这是定理，这是定理，这是定理，这是定理，这是定理，这是定理，这是定理，这是定理，这是定理，这是定理，这是定理，这是定理，这是定理，这是定理，这是定理，这是定理，这是定理，这是定理，这是定理，这是定理，这是定理，这是定理，这是定理，这是定理，这是定理，这是定理。
\end{theorem}
\begin{proof}
  这是证明，这是证明，这是证明，这是证明，这是证明，这是证明，这是证明，这是证明，这是证明，这是证明，这是证明，这是证明，这是证明，这是证明，这是证明，这是证明，这是证明，这是证明，这是证明，这是证明，这是证明，这是证明。
\end{proof}
\begin{definition}
  这是定义，这是定义，这是定义，这是定义，这是定义，这是定义，这是定义，这是定义，这是定义，这是定义，这是定义，这是定义，这是定义，这是定义，这是定义，这是定义，这是定义，这是定义，这是定义，这是定义，这是定义，这是定义，这是定义，这是定义，这是定义
\end{definition}
\begin{lemma}
  这是引理
\end{lemma}
现有成果介绍, 现有成果介绍, 现有成果介绍, 现有成果介绍, 现有成果介绍.
现有成果介绍, 现有成果介绍, 现有成果介绍, 现有成果介绍, 现有成果介绍.
现有成果介绍, 现有成果介绍, 现有成果介绍, 现有成果介绍, 现有成果介绍.

\section {本研究工作内容和章节安排}


\chapter {相关技术原理的分析与研究}


\section {xxx相关技术原理分析与研究}


\section {本章小节}


这是小节，这是小节，这是小节，这是小节，这是小节，这是小节，
这是小节，这是小节，这是小节，这是小节，这是小节，这是小节，
这是小节，这是小节，这是小节，这是小节，这是小节，这是小节，
这是小节，这是小节，这是小节，这是小节，这是小节，这是小节，

这是小节，这是小节，这是小节，这是小节，这是小节，这是小节，
这是小节，这是小节，这是小节，这是小节，这是小节，这是小节，
这是小节，这是小节，这是小节，这是小节，这是小节，这是小节，
这是小节，这是小节，这是小节，这是小节，这是小节，这是小节，

\clearpage %{\pagestyle{empty}\cleardoublepage}

\chapter {相关算法的研究}


\section {相关技术原理简述}


\subsection {xxx算法简述}\label{section:acoustic-feat}

\subsection {xxx模型简述}



\begin{equation}
\label{eq:DFT}
A_k = \sum_{n=0}^{N-1} {e^{-i {2 \pi \over N}}}^{kn} a_n 
\end{equation}


这是对公式的引用,\autoref{eq:DFT}中的引用

\section{本章小节}

本章我们对设计的xxx架构进行了详细的分析与实际性能对比测试。



\chapter {总结与展望}

\section {工作总结}

这是总结, 这是总结, 这是总结, 这是总结, 这是总结, 
这是总结, 这是总结, 这是总结, 这是总结, 这是总结, 
这是总结, 这是总结, 这是总结, 这是总结, 这是总结, 
这是总结, 这是总结, 这是总结, 这是总结, 这是总结, 

\newpage
\section {工作展望}

这是展望, 这是展望, 这是展望, 这是展望, 这是展望.
这是展望, 这是展望, 这是展望, 这是展望, 这是展望.
这是展望, 这是展望, 这是展望, 这是展望, 这是展望.
这是展望, 这是展望, 这是展望, 这是展望, 这是展望.

\clearpage%{\pagestyle{empty}\cleardoublepage}
\backmatter
\linespread{1.6}\selectfont
%%%% ===== 参考文献 =====
\setlength{\bibsep}{1ex}  % 需 natbib 宏包
\begin{thebibliography}{99}
\addcontentsline{toc}{chapter}{参考文献}
\thispagestyle{fancy}

% \addtolength{\itemsep}{-5pt}


\bibitem{bi:hw-ml-challenges}
Sze V, Chen Y H, Emer J, et al. Hardware for machine learning: Challenges and opportunities[C]//2017 IEEE Custom Integrated Circuits Conference (CICC). IEEE, 2017: 1-8.

\bibitem{bi:hw-sw-co-design}
Brian Bailey, Hardware-Software Co-Design Reappears. semiengineering .2019-07-25. https://semiengineering.com/hardware-software-co-design-reappears/
 



\end{thebibliography}


\clearpage{\pagestyle{fancy}\cleardoublepage}
\chapter*{附录}
\addcontentsline{toc}{chapter}{附录}
附录部分, 附录部分, 附录部分, 附录部分, 附录部分,
附录部分, 附录部分, 附录部分, 附录部分, 附录部分,
附录部分, 附录部分, 附录部分, 附录部分, 附录部分,
附录部分, 附录部分, 附录部分, 附录部分, 附录部分.

附录部分, 附录部分, 附录部分, 附录部分, 附录部分,
附录部分, 附录部分, 附录部分, 附录部分, 附录部分,
附录部分, 附录部分, 附录部分, 附录部分, 附录部分,
附录部分, 附录部分, 附录部分, 附录部分, 附录部分.

\newpage 

附录部分, 附录部分, 附录部分, 附录部分, 附录部分,
附录部分, 附录部分, 附录部分, 附录部分, 附录部分,
附录部分, 附录部分, 附录部分, 附录部分, 附录部分,
附录部分, 附录部分, 附录部分, 附录部分, 附录部分.






\clearpage{\pagestyle{fancy}\cleardoublepage}
\chapter*{致谢}
\addcontentsline{toc}{chapter}{致谢}

\linespread{1.4}\selectfont


“日月忽其不淹兮，春与秋其代序”，丽娃河畔，共青路边，梧桐落木，一如既往。两年前，我乘着开往这座城市的列车，来到华师大，努力想要看清这周围的每一个角落。千来个日夜，莘莘研学之路，终以这论文的结束而告一段落。


如今，这两年多的时光还可算明了。人生路漫漫，但难免有所遗憾。正如这篇毕业论文，撰写接近尾声之际，仍然还有许多尚待修补的地方，限于当下能力有限，只好尽力而为。也许期年之后，回顾此文，自惭之至，但也可看清来时之路。



\begin{flushright}
张三

于数学科学学院

2023年3月
\end{flushright}



\iffalse  %科研成果，貌似不用有需要的加上。

\clearpage{\pagestyle{fancy}\cleardoublepage}
\chapter*{研究成果}
\addcontentsline{toc}{chapter}{研究成果}

软件蓍作权：基于xxxx系统 V1.0

\fi

\end{document}
